% Materialized canonical-collar bridge and conditional arbitrary-field reduced proofs.
% Date: 28 August 2026.  Candidate only; bridge production verification pending.

% Proof of the canonical-collar repair of the source-legal arbitrary-field
% virtual-annulus comparison.
% Date: 28 August 2026.
% Standalone production-verification candidate; not integrated into the book.
% The rejected c5a643/292bb and 0315d3/c903cf artifacts are preserved
% unchanged in separate files.

\begin{proof}[Proof of
\cref{lem:p21p3-gf-legal-source-localization}]
Use the ordinary local moment description in
\cite[Proposition~15.1.1(b)]{CLD25}.  After shrinking \(W\) about \(y\),
there are an analytic moment \(F_{\mathrm{form}}\) and a smooth superform
\(\eta\) such that \(\alpha=F_{\mathrm{form}}^*\eta\).  Shrink once more so that
the prescribed finite block \(G\) is defined.  When \(\chi\) is present,
\cite[Corollary~15.1.2]{CLD25} gives, after one further shrinking, an analytic
moment \(H\) and a smooth function \(\Phi\) such that
\(\chi=\Phi\circ H_{\mathrm{trop}}\).  Relabel this common neighborhood as
\(W\).

Apply \cite[Lemma~12.3.6]{CLD25} to the analytic space \(W\) at \(y\).  It
gives a compact analytic neighborhood \(V\) of \(y\) in \(W\) and a moment
\(F_{\mathrm{bd}}\) such that
\[
 F_{\mathrm{bd,trop}}(y)\notin
 F_{\mathrm{bd,trop}}(\partial(V/W)).
\]
Concatenate \(F_{\mathrm{bd}}\), \(F_{\mathrm{form}}\), \(G\), and \(H\).
Projection to the first block proves
\eqref{eq:p21p3-gf-source-local-boundary}; projection to the form and cutoff
blocks proves \eqref{eq:p21p3-gf-source-local-form} and
\eqref{eq:p21p3-gf-source-cutoff-factorization}.  Equality of two full
tropical values implies equality of their \(H\)-projections, so applying
\(\Phi\) proves \eqref{eq:p21p3-gf-source-cutoff-fibre-constancy}.

The remainder is the localization used in the proof of
\cite[Theorem~16.8.5]{CLD25}.  The set \(\Omega\) in
\eqref{eq:p21p3-gf-source-local-omega} is open because the relative boundary
is compact.  The set \(U\) misses \(\partial(V/W)\), hence is contained in the
relative interior of \(V\) and is open in \(W\).  It contains \(y\) by
\eqref{eq:p21p3-gf-source-local-boundary}.  Moreover
\(F_{\mathrm{trop}}|_U\) is obtained from the proper map
\(F_{\mathrm{trop}}:V\to\mathbf R^N\) by base change along
\(\Omega\hookrightarrow\mathbf R^N\), and is therefore proper.  Thus the
inverse image of the compact support of \(\lambda\) is compact in \(U\).
The displayed pullback form consequently extends by zero across
\(W\setminus U\), and the equality with \(\alpha\) where \(\lambda=1\)
follows from \eqref{eq:p21p3-gf-source-local-form}.

For the last assertion, apply the construction at each point of the compact
test support, starting inside the prescribed member of the source cover.
The resulting open sets \(U\) cover that support, so finitely many suffice.
Smooth cutoffs subordinate to this finite cover, with sum one near the
compact support, exist by \cite[Corollary~15.2.9]{CLD25}.  Current identities
are local and therefore glue after this finite decomposition.  If one of
those cutoffs is retained in a
single local presentation, the preceding application of
\cite[Corollary~15.1.2]{CLD25} supplies \(H\); at no point was the cutoff
itself used as an analytic coordinate.
\end{proof}

\begin{proof}[Proof of
\cref{prop:p21p3-gf-virtual-annulus-bridge}]
The residual curve and the branch--valuation correspondence over an
arbitrary complete field are
\cite[\S\S3.3.5.2, 4.2.6, 4.2.9, and 4.2.11.1]
{Ducros24CurveStructure}.  Since \(x\)
is interior, every closed point of \(C_x\) occurs as a branch.  For the branch
\(b\) corresponding to \(p_b\), the refined value group is the lexicographic
extension of \(|\mathcal H(x)^{\times}|\) by the discrete valuation
\(\operatorname{ord}_{p_b}\).  If \(s_b\) is a uniformizer of the discrete
valuation ring \(\mathcal O_{C_x,p_b}\), then
\cite[\S4.2.11.1]{Ducros24CurveStructure} says precisely that the class of
its refined value generates
\[
 |\mathcal H(x)^{\times}|_b/|\mathcal H(x)^{\times}|\simeq\mathbf Z.
\]
Since \(Y\) is normal and \(x\) is nonrigid,
\cite[\S4.2.18]{Ducros24CurveStructure} identifies
\(\mathcal O_{Y,x}\) with \(\kappa(x)\), and
\cite[\S4.2.6]{Ducros24CurveStructure} says that \(\kappa(x)\) is dense in
\(\mathcal H(x)\).  Hence \(s_b\), viewed as a homogeneous residue element
of degree one, has a norm-one lift \(f_b\in\mathcal O_{Y,x}\).  This is also
the finite residue-lifting step made explicit in the proof of
\cite[Lemma~6.14]{GJR26}.  Its value in \(\mathcal H(x)\) is nonzero, so
after shrinking it is an analytic unit.  This proves that
\eqref{eq:p21p3-gf-branch-order-definition} is meaningful.  Replacing a
uniformizer by another one multiplies it by a unit in
\(\mathcal O_{C_x,p_b}\), whose class in the displayed quotient is trivial.
This proves independence and also gives
\eqref{eq:p21p3-gf-branch-order-residue}.

We next prove the exact ray assertion.  By
\cite[Proposition~4.3.9(2) and \S\S4.3.10.2--4.3.10.3]
{Ducros24CurveStructure}, \(b\) has a cofinal family of virtual-annular
sections.  Choose one and call it \(Z\).  Its analytic skeleton and canonical
retraction are compatible with every complete scalar extension by
\cite[\S\S3.6.35.5--3.6.35.10]{Ducros24CurveStructure}.  After a finite
separable extension and then passage to a completed algebraic closure, a
component above \(Z\) is an ordinary annulus whose skeleton maps
homeomorphically to \(S^{\mathrm{an}}(Z)\).

Let \(\beta\) be the branch of this split annulus above \(b\), and choose an
annulus coordinate \(\theta\) so that \(-\log|\theta|\) increases away from
the endpoint over \(x\).
The annular calculation in
\cite[\S\S3.6.22--3.6.24 and \S4.3.6.4]
{Ducros24CurveStructure} shows that, after shortening the annulus once for
the finite family \(g_1,\ldots,g_m\),
\begin{equation}
\label{eq:p21p3-gf-split-annulus-monomials}
                  |g_j|=c_j|\theta|^{n_j}
                  \qquad(1\leq j\leq m)
\end{equation}
for constants \(c_j>0\) and integers \(n_j\).  Naturality of the refined
branch valuation is \cite[\S4.2.9]{Ducros24CurveStructure}.  Consequently,
restriction from \(b\) to \(\beta\) sends the generator defined by \(g_m\) to
a positive integral multiple of the primitive annular generator.  Hence
there is one integer \(e_b>0\), independent of \(j\), such that
\begin{equation}
\label{eq:p21p3-gf-common-expansion-index}
                       n_j=e_b\partial_b(g_j)
                       \qquad(1\leq j\leq m).
\end{equation}
Indeed, this is just the homomorphism between the two rank-one quotient
value groups; \(n_m=e_b\) because \(\partial_b(g_m)=1\).  The same integer
is the tangent-lattice expansion index of this split annular presentation
by \cite[\S3.6.24.2]{Ducros24CurveStructure}.  In the finite-morphism
setting of \cite[Theorem~4.3.13(i)]{Ducros24CurveStructure}, that annular
index is identified with the corresponding residual ramification index.

Put
\[
 t=-\log|g_m|+\log|g_m(x)|.
\]
The orientation and
\cite[\S4.2.11.1]{Ducros24CurveStructure} make \(t\) strictly positive
away from \(x\) on a sufficiently short section.  Equations
\eqref{eq:p21p3-gf-split-annulus-monomials} and
\eqref{eq:p21p3-gf-common-expansion-index}, followed by passage to the
limit at \(x\), give on the skeleton
\begin{equation}
\label{eq:p21p3-gf-coordinate-ray-formula}
 -\log|g_j|+\log|g_j(x)|
       =\partial_b(g_j)t
       \qquad(1\leq j\leq m).
\end{equation}
Every invertible analytic function on a virtual annulus has the same norm
as at its canonical retraction by
\cite[\S3.6.35.8]{Ducros24CurveStructure}.  Thus
\eqref{eq:p21p3-gf-coordinate-ray-formula} holds on the whole shortened
section \(Z_b\), and \(t\) identifies its skeleton with an interval
\((0,\varepsilon_b)\).  This proves
\eqref{eq:p21p3-gf-exact-branch-ray}.  Since the last coordinate of \(v_b\)
is one, \(v_b\) is primitive in the saturated ambient lattice.  The
one-dimensional skeleton attached to the moment, restricted to any compact
subannular band, has the unique integral affine structure for which the
moment maps its faces integral-affinely onto their images by
\cite[\S8.13]{GJR26}; hence this is also the primitive tangent of the
corresponding source one-cell.  This proves
\cref{itm:p21p3-gf-bridge:ray}.

We prove the finite separation assertion.  Let \(p_i=p_{b_i}\).  Weak
approximation on the normal projective curve \(C_x\) gives rational functions
\(s_1,\ldots,s_s\) with
\begin{equation}
\label{eq:p21p3-gf-kronecker-orders}
                \operatorname{ord}_{p_i}(s_j)=\delta_{ij}.
\end{equation}
For completeness, choose a closed point \(q\) away from the \(p_i\).  For
\(n\gg0\), Serre vanishing gives
\[
 H^1\!\left(C_x,
   \mathcal O_{C_x}\!\left(nq-2\sum_i p_i\right)\right)=0.
\]
The divisor exact sequence makes the simultaneous value-and-first-jet map
from \(H^0(C_x,\mathcal O_{C_x}(nq))\) onto
\(\bigoplus_i\mathcal O_{C_x,p_i}/\mathfrak m_{p_i}^2\) surjective.
Prescribing a nonzero first jet and value zero at \(p_j\), and value one at
the other \(p_i\), yields \(s_j\) satisfying
\eqref{eq:p21p3-gf-kronecker-orders}.  This argument is valid for closed
points of arbitrary residue degree.

Lift the \(s_j\) to norm-one analytic units by the density argument above
(equivalently, the lifting step in the proof of
\cite[Lemma~6.14]{GJR26}), and first append those lifts to \(F_0\).  By
\cite[Proposition~4.2.4(ii)]{Ducros24CurveStructure}, every
analytic unit has nonzero branch order in only finitely many branches.
The preliminary outgoing vector in branch \(b_i\) therefore has, in the new
\(s_1,\ldots,s_s\) block, the standard basis vector \(e_i\); in particular it
is primitive, and the selected vectors are pairwise nonparallel.

There can still be an unselected branch whose preliminary vector is a
positive multiple of the vector for some \(b_i\).  There are only finitely
many such branches: any one of them has positive \(s_i\)-order and hence
belongs to the finite exceptional set supplied by
\cite[Proposition~4.2.4(ii)]{Ducros24CurveStructure}.  Call their union
\(Q\).  For every \(c\in Q\), weak approximation as above gives a rational
function \(h_c\) satisfying
\begin{equation}
\label{eq:p21p3-gf-collision-killers}
 \operatorname{ord}_{p_{b_i}}(h_c)=0\quad(1\leq i\leq s),
 \qquad
 \operatorname{ord}_{p_c}(h_c)=1.
\end{equation}
Lift all \(h_c\) and append them as well.  The new coordinates vanish in
branch order on every selected branch, whereas the \(h_c\)-coordinate has
order one on \(c\).  Thus each old collision is destroyed.  A newly active
branch cannot create a new collision: projection to the preliminary
coordinates would already have made it an element of \(Q\).  This proves
\cref{itm:p21p3-gf-bridge:separation}.

We finally prove \cref{itm:p21p3-gf-bridge:localization}.  First use
\cref{lem:p21p3-gf-legal-source-localization} inside the prescribed source
\(W\), with no branch block yet, and retain the resulting moment germ
\(F_0\).  Apply
\cref{itm:p21p3-gf-bridge:separation} to \(F_0\).  On a common neighborhood
of \(x\), let \(F_{\mathrm{br}}\) be the resulting separating block and put
\(G_0=(F_0,F_{\mathrm{br}})\).  Let \(O\subset W\) be a common open
neighborhood of \(x\) on which every coordinate of \(G_0\) is an analytic
unit.  We choose the source star before adding any further moment coordinates.
By \cite[Theorem~4.5.4]{Ducros24CurveStructure}, after shrinking inside
\(O\), there is a connected open elementary neighborhood \(V_*\) of \(x\)
whose closure is a compact tree and for which
\begin{equation}
 V_*\setminus\{x\}
   =\coprod_{\lambda\in\Lambda}B_\lambda
      \amalg\coprod_{\nu=1}^{r}A_\nu .
 \label{eq:p21p3-gf-elementary-star-decomposition}
\end{equation}
The \(B_\lambda\) are virtual discs, whereas
\(A_1,\ldots,A_r\) are the finitely many gently virtual-annular components.
Every component in \eqref{eq:p21p3-gf-elementary-star-decomposition}
represents one branch at \(x\), and the compact-tree condition prevents an
annular component from returning to \(x\) at its other end.

An invertible analytic function has constant norm on a virtual disc by
\cite[\S3.6.32.3]{Ducros24CurveStructure}.  Continuity at \(x\) therefore
gives
\begin{equation}
 G_{0,\mathrm{trop}}|_{B_\lambda}
       =G_{0,\mathrm{trop}}(x)
       \qquad(\lambda\in\Lambda).
 \label{eq:p21p3-gf-disc-components-at-vertex}
\end{equation}
This retains all, possibly infinitely many, virtual-disc components whole.

Only the finitely many annular components need to be shortened.  Let
\(b_\nu\) be the branch represented by \(A_\nu\), and put
\(w_\nu=(\partial_{b_\nu}(g))_g\), where \(g\) runs through the coordinates
of \(G_0\).  For this calculation only, choose a norm-one analytic unit
\(u_\nu\in\mathcal O_{Y,x}^{\times}\) with
\(\partial_{b_\nu}(u_\nu)=1\); it is not appended to the moment.  Since there
are only finitely many \(A_\nu\), all the \(u_\nu\) are units on a common
sufficiently short collar.  After a finite separable extension splitting
\(A_\nu\), choose an annulus coordinate \(\theta_\nu\) oriented away from the
end over \(x\).  Apply \eqref{eq:p21p3-gf-split-annulus-monomials}
simultaneously to the coordinates of \(G_0\) and to \(u_\nu\).  Naturality of
the refined branch valuation, \cite[\S4.2.9]{Ducros24CurveStructure}, and the
rank-one quotient argument of
\eqref{eq:p21p3-gf-common-expansion-index} give one integer \(e_\nu>0\) such
that the split-annulus exponents are
\[
 n_{\nu,g}=e_\nu\partial_{b_\nu}(g),
 \qquad n_{\nu,u}=e_\nu.
\]
Consequently, after choosing a skeleton point \(q_\nu\) sufficiently close to
the end accumulating at \(x\),
\begin{equation}
 G_{0,\mathrm{trop}}(z)
   =G_{0,\mathrm{trop}}(x)+t_\nu(z)w_\nu,
 \qquad 0<t_\nu(z)<\delta_\nu,
 \label{eq:p21p3-gf-finite-annulus-ray-formula}
\end{equation}
on the retained collar, where
\[
 t_\nu(z)=-\log|u_\nu(z)|+\log|u_\nu(x)|.
\]
The orientation makes \(t_\nu\) positive and strictly increasing along the
retained skeleton.  By \cite[\S3.6.35.8]{Ducros24CurveStructure} it factors
through the canonical retraction.  If \(w_\nu=0\), the displayed formula says
that the whole retained collar maps to the vertex.
Sections~3.6.35.9--3.6.35.10 of \cite{Ducros24CurveStructure} identify the
inverse image \(A_\nu^0\) of the retained skeleton interval with an analytic
virtual subannulus, compatibly with scalar extension.

Let \(W_1\) be the component containing \(x\) of
\(V_*\setminus\{q_1,\ldots,q_r\}\).  Analytic curves are locally connected
graphs by \cite[Theorem~3.5.1]{Ducros24CurveStructure}; hence \(W_1\) is an
open connected analytic neighborhood of \(x\).  The tree description and
the choice of the \(q_\nu\) give
\begin{equation}
 W_1\setminus\{x\}
   =\coprod_{\lambda\in\Lambda}B_\lambda
      \amalg\coprod_{\nu=1}^{r}A_\nu^0.
 \label{eq:p21p3-gf-trimmed-star-decomposition}
\end{equation}
Thus no independent truncation has been made in infinitely many directions.

Equation \eqref{eq:p21p3-gf-disc-components-at-vertex} and the order-one
separator coordinate show that a selected branch cannot be represented by a
virtual-disc component.  Hence it is represented by a unique
\(A_{\nu(i)}^0\).  Put \(w_i=w_{\nu(i)}\) and
\(\delta_i=\delta_{\nu(i)}\).  The residue-unit block contains a coordinate
of order one on \(b_i\), so \(w_i\ne0\), and
\cref{itm:p21p3-gf-bridge:separation} says that
\(\mathbf R_{>0}w_i\) is different from the ray of every other active
component.  After moving the finitely many \(q_\nu\) closer to \(x\), if
necessary, equations \eqref{eq:p21p3-gf-disc-components-at-vertex} and
\eqref{eq:p21p3-gf-finite-annulus-ray-formula} therefore give the exact
preliminary locator
\begin{equation}
 G_{0,\mathrm{trop}}^{-1}\!\left(
 G_{0,\mathrm{trop}}(x)+(0,\delta_i)w_i\right)\cap W_1
 = A_{\nu(i)}^0.
\end{equation}
Indeed, every virtual disc and every inactive annulus maps to the vertex,
every other active annulus maps into a distinct open ray, and on the selected
annulus the order-one coordinate makes \(t_\nu\) a genuine radial parameter.
This excludes remote source collisions without an infinite simultaneous
choice of branch sections.

Reapply \cref{lem:p21p3-gf-legal-source-localization} inside \(W_1\), with
\(G=G_0\), and call its compact domain and full moment \(V\) and \(F\).  The
lemma gives boundary avoidance for this full moment.  Hence the open set
\(U\) in \eqref{eq:p21p3-gf-source-local-omega} contains \(x\), lies in
\(W_1\), and \(F_{\mathrm{trop}}|_U\) is proper.  Let
\(P:\mathbf R^N\to\mathbf R^{N_0}\) be projection to the \(G_0\)-block, so
that \(P\circ F_{\mathrm{trop}}=G_{0,\mathrm{trop}}\).

For a selected branch \(b_i\), the full moment \(F\) still contains the
order-one separator coordinate.  On \(A_{\nu(i)}^0\), the formula above for
that coordinate identifies its normalized negative logarithm with the radial
parameter \(t_i=t_{\nu(i)}\).  Since \(U\) is an open neighborhood of \(x\),
\cite[\S4.3.10.3]{Ducros24CurveStructure} and
\cite[\S\S3.6.35.9--3.6.35.10]{Ducros24CurveStructure} allow us to choose
\(0<\varepsilon_i<\delta_i\) so that
\begin{equation}
 Z_{b_i}:=t_i^{-1}((0,\varepsilon_i))
       \subset A_{\nu(i)}^0\cap U
 \label{eq:p21p3-gf-canonical-full-source-collar}
\end{equation}
is the full canonical virtual subannulus over that skeleton interval.  After
decreasing \(\varepsilon_i\) once more, the calculation in
\cref{itm:p21p3-gf-bridge:ray}, applied simultaneously to all coordinates of
\(F\), gives
\[
 F_{\mathrm{trop}}(z)=F_{\mathrm{trop}}(x)+t_i(z)v_i
 \qquad(z\in Z_{b_i}),
\]
where \(v_i\) is the full primitive branch-order vector.  This proves
\eqref{eq:p21p3-gf-source-separated-ray-image}; it also gives
\(P(v_i)=w_i\).  If \(z\in U\) and
\(F_{\mathrm{trop}}(z)=F_{\mathrm{trop}}(x)+t v_i\) for
\(0<t<\varepsilon_i\), projection and
\eqref{eq:p21p3-gf-preliminary-exact-source-locator} put \(z\) in
\(A_{\nu(i)}^0\).  The order-one separator coordinate then gives
\(t_i(z)=t\).  By the definition
\eqref{eq:p21p3-gf-canonical-full-source-collar}, \(z\in Z_{b_i}\).  The
reverse inclusion follows from the displayed ray formula, and hence
\eqref{eq:p21p3-gf-source-separated-ray-preimage} holds.

Consequently the Euclidean current
\[
 \langle T_{\Omega},\beta\rangle
       =\langle T|_U,F^*\beta\rangle
\]
used in the proof of \cite[Theorem~16.8.5]{CLD25}, when tested on a
sufficiently short relative-open subsegment of the selected ray, receives
source contribution only from \(Z_{b_i}\).  If a norm-one lift of
\(g\in\operatorname{Rat}(C_x)^\times\) is appended, projection to the
previous coordinates preserves the exact separation, while
\eqref{eq:p21p3-gf-branch-order-residue} identifies its scalar slope with
\(\operatorname{ord}_{p_{b_i}}(g)\).  Multiplication by the edge flux \(a_i\)
gives the asserted balancing term.  No smooth plateau was used as a moment
coordinate, and no scalar extension of \(T\) occurred.

Equation \eqref{eq:p21p3-gf-common-expansion-index} also proves the claim in
\cref{rem:p21p3-gf-virtual-annulus-index}.  Under pushforward, the rank-one
matrix current has one factor \(e_b\): its tangent tensor contributes
\(e_b^2\) and one-dimensional change of measure contributes \(e_b^{-1}\).
The canonical calibration has the same expansion factor by
\cite[\S14.3.2, Lemma~14.3.4, and Proposition~14.4.2]{CLD25}.  No arbitrary
current was extended to a splitting field at any point of the argument.
\end{proof}

% Source-aware conditional proof for
% p21p3_curve_rigidity_arbitrary_field_canonical_collar_conditional_statement_20260828.tex.
% Date: 28 August 2026.
% The only conditional inputs are the canonical-collar source-localization
% lemma and prop:p21p3-gf-virtual-annulus-bridge at the exact hashes stated
% in the accompanying conditional statement.

\paragraph{Imported facts and the conditional input.}
Currents, positivity, multiplication of an order-zero current by a
continuous function, and differentials of currents are used as in
\cite[\S\S7.1--7.2, Example~7.2.12, and \S7.7]{CLD25}.  In particular,
\cite[Lemma~7.7.5]{CLD25} gives
\begin{equation}
\label{eq:p21p3-gfc-transpose-dd}
 \bigl\langle \mathrm d'\mathrm d''S,\varphi\bigr\rangle
 =
 \bigl\langle S,\mathrm d'\mathrm d''\varphi\bigr\rangle .
\end{equation}
Positive currents of bidimension $(0,0)$ are Radon measures by
\cite[Proposition~7.2.8]{CLD25}, and positive currents in every bidegree
have order zero by \cite[Proposition~7.2.9]{CLD25}.  We use the local
moment description of smooth forms and functions from
\cite[Proposition~15.1.1(b) and Corollary~15.1.2]{CLD25}, smooth
partitions of unity and subordinate cutoffs from
\cite[Corollary~15.2.4, Proposition~15.2.8, and
Corollary~15.2.9]{CLD25}, boundary-avoiding moments from
\cite[Lemma~12.3.6]{CLD25}, and the proper tropical pushforward in the
proof of \cite[Theorem~16.8.5]{CLD25}.  We also use
\cite[Proposition~3.3.6, Lemma~12.3.6, Theorem~12.2.4,
Corollary~12.2.5, \S14.3.2, Lemma~14.3.4,
Proposition~14.4.2, \S16.6, and Proposition~16.6.3]{CLD25}.

For an arbitrary complete ground field, the curve structure used below is
the branch--valuation correspondence, the virtual-annulus description,
and triangulation theory in
\cite[\S\S3.3.5.2, 4.2.6, 4.2.9, 4.2.10--4.2.11, and
4.3.9--4.3.10]{Ducros24CurveStructure}.  Thus an interior
type-$2$ point $x$ has a normal projective integral residual curve
$C_x$, and its branches correspond to all closed points of $C_x$.
The branch corresponding to $p$ carries the discrete valuation
$\operatorname{ord}_p$.  An interior type-$3$ point has two branches
and a virtual-annular neighborhood.  For the arbitrary-field endpoint
statement used later, \cite[Theorem~4.2.5]{Ducros24CurveStructure}
shows that an interior nonrigid point with two branches must have type
$2$ or $3$, while \cite[Theorem~4.5.4, especially \S4.5.4.1, and
\S4.5.19.1]{Ducros24CurveStructure} supplies the virtual-disc local
model and confirms that type-$1$ and type-$4$ points of a rig-smooth
curve are unibranch.  The finite-fibre and integral-affine
properties of tropical skeleta are
\cite[Corollary~3.18 and Remark~3.20]{GJR26}; finite families of
homogeneous residue functions are lifted to analytic germs by the
construction in the proof of \cite[Lemma~6.14]{GJR26}.

The proof below assumes the canonical-collar candidate
\cref{lem:p21p3-gf-legal-source-localization} and exactly the three
assertions of the canonical-collar candidate
\cref{prop:p21p3-gf-virtual-annulus-bridge}.  In particular, it uses the
exact downstairs branch ray, finite separation of prescribed type-$2$
branches, the canonical-collar elementary star in
\cref{itm:p21p3-gf-bridge:localization}, its
global-in-\(U\) exact-preimage identity, and compatibility with the
source-separated localization of \cite[Theorem~16.8.5]{CLD25}.  It never
pulls an arbitrary current to a field extension.  Nothing in this proof is
presented as established until the conditional input passes independent and
production verification.

\begin{proof}[Proof of \cref{lem:p21p3-gfc-rigid-zero-atom}]
The regular one-dimensional local ring $\mathcal O_{Y,y}$ has a
uniformizer represented, after shrinking, by an analytic function $f$
with $f(y)=0$.  Regard $f$ as a morphism to
$\mathbf A^{1,\mathrm{an}}_K$.  Its fibre germ at $y$ is
zero-dimensional.  Since source and target are boundaryless, the
boundary-composition identity \cite[(11.4.1.3)]{CLD25} gives
$y\in\operatorname{Int}(Y/\mathbf A^{1,\mathrm{an}}_K)$, and
\cite[Lemma~11.6.5(a)]{CLD25} makes $f$ finite at $y$.  The induced
finite homomorphism from the discrete valuation ring at $0$ to
$\mathcal O_{Y,y}$ is torsion-free and hence flat.  After shrinking as
in \cite[\S11.6.3]{CLD25}, choose a compact analytic domain $Q_{R_0}$
containing $y$, an admissible radius $R_0>0$, and a finite flat map of
constant degree $d>0$
\begin{equation}
\label{eq:p21p3-gfc-finite-rigid-tube}
             f:Q_{R_0}\longrightarrow D(0,R_0)
\end{equation}
such that $f^{-1}(0)=\{y\}$.  A finite map has empty relative boundary;
the same boundary identity puts the relative boundary of $Q_{R_0}$
over $\{|T|=R_0\}$.

Choose an admissible radius $0<R<R_0$ and fix
\[
 0<\varepsilon_0<\tfrac12\log(R_0/R).
\]
For an admissible $0<r<R$ with
$\log(R/r)>4\varepsilon_0$, put $s_R=-\log R$ and
$s_r=-\log r$.  Choose a $C^\infty$ function
$\rho_r:\mathbf R\to[0,1]$ which is zero on
\[
 (-\infty,s_R-\varepsilon_0]\cup[s_r+\varepsilon_0,\infty),
\]
is one on $[s_R,s_r]$, is nondecreasing on
$[s_R-\varepsilon_0,s_R]$, and is nonincreasing on
$[s_r,s_r+\varepsilon_0]$.  Set
\[
 q_r(t):=\int_{-\infty}^{t}\rho_r(\tau)\,\mathrm d\tau,
 \qquad
 H_r:=\int_{\mathbf R}\rho_r(\tau)\,\mathrm d\tau .
\]
Then
\begin{equation}
\label{eq:p21p3-gfc-radial-data}
 H_r\geq\log(R/r),
 \qquad
 \int_{\mathbf R}|q_r''(t)|\,\mathrm dt=2.
\end{equation}
On the punctured tube put
$\chi_r=q_r(-\log|f|)$.  Extend it by zero across the outer region and
by the constant $H_r$ near $y$.  The extensions occur where $q_r$
is constant, so $\chi_r$ is a compactly supported CLD-smooth function,
with support in the fixed compact domain $Q_{R_0}$.

On each of the two transition tubes, the map $s=-\log|f|$ projects onto
the corresponding target interval and
\begin{equation}
\label{eq:p21p3-gfc-radial-dd}
 \mathrm d'\mathrm d''\chi_r
   =q_r''(s)\,\mathrm d's\wedge\mathrm d''s .
\end{equation}
The sign of $q_r''$ is fixed on each transition band.  Restrict proper
calibrated pushforward under
\eqref{eq:p21p3-gfc-finite-rigid-tube} to either transition tube.  By
\cite[Proposition~14.4.2]{CLD25} it is the corresponding target annulus
with total local degree $d$; in particular all sheets and residue-degree
factors are included.  If $\mu_f$ and $\mu_D$ denote the source and radial
calibrated measures, then $f_*\mu_f=d\mu_D$ on each transition band.
Since the two bands have disjoint supports and $q_r''$ has one sign on
each of them, \eqref{eq:p21p3-gfc-radial-data} gives
\begin{equation}
\label{eq:p21p3-gfc-radial-tv}
 \left\|\mathrm d'\mathrm d''\chi_r\wedge\delta_Y\right\|_{\mathrm{TV}}
 =d\int_{\mathbf R}|q_r''(t)|\,\mathrm dt=2d.
\end{equation}

The measure $\mu_u=\mathrm d'\mathrm d''[u]$ is positive by
\cite[Lemma~8.5.5]{CLD25}.  Positivity,
\eqref{eq:p21p3-gfc-transpose-dd}, and
\eqref{eq:p21p3-gfc-radial-tv} yield
\begin{align}
 \log(R/r)\,\mu_u(\{y\})
 &\leq \int_Y\chi_r\,\mathrm d\mu_u                                      \notag\\
 &=\bigl\langle[u],\mathrm d'\mathrm d''\chi_r\bigr\rangle                \notag\\
 &\leq 2d\sup_{Q_{R_0}}|u|.
\label{eq:p21p3-gfc-radial-bound}
\end{align}
The right side is independent of $r$.  Nontriviality of the valuation
supplies admissible radii tending to zero, so
\eqref{eq:p21p3-gfc-radial-bound} proves $\mu_u(\{y\})=0$.
\end{proof}

\begin{proof}[Proof of \cref{prop:p21p3-gfc-smooth-rigidity}]
We first record three ingredients and then globalize them.

\paragraph{The Euclidean calculation.}
Let $\Gamma$ be a locally finite one-dimensional polyhedral set in an
open subset of $\mathbf R^N$, and let $S$ be a positive current of
bidimension $(1,1)$, supported on $\Gamma$, with
$\mathrm d'S=\mathrm d''S=0$.  Positivity and
\cite[Proposition~7.2.9]{CLD25} write $S$ as a symmetric positive
semidefinite matrix $M=(M_{jk})$ of Radon measures.  Closedness is the
distributional system
\begin{equation}
\label{eq:p21p3-gfc-matrix-divergence}
                   \sum_j\partial_jM_{jk}=0
                   \quad(1\leq k\leq N).
\end{equation}
On the relative interior of an edge, choose affine coordinates
$(t,z_2,\ldots,z_N)$ in which the edge is
$\{z_2=\cdots=z_N=0\}$.  Testing
\eqref{eq:p21p3-gfc-matrix-divergence} against
$\psi(t)z_\ell$, for $\ell>1$, kills every transverse row and column;
testing against $\psi(t)$ then makes the remaining measure a constant
multiple of Lebesgue measure.  Thus
\begin{equation}
\label{eq:p21p3-gfc-edge-density}
                     S|_e=b_e[\,e\,],
                     \qquad b_e\geq0,
\end{equation}
with $b_e$ constant.  Endpoint terms at a vertex give
\begin{equation}
\label{eq:p21p3-gfc-balancing}
                       \sum_{e\ni v}b_eu_{e,v}=0,
\end{equation}
where $u_{e,v}$ is the primitive outgoing tangent.  A positive remainder
at $v$ would have the form $P\delta_v$ with $P$ positive
semidefinite.  Once the endpoint terms balance, tests vanishing at $v$
with arbitrary first jet isolate the derivatives of $P\delta_v$ in
\eqref{eq:p21p3-gfc-matrix-divergence}; hence $P=0$.  In the vertex
applications below,
\cref{itm:p21p3-gf-bridge:localization} produces a one-vertex source
neighborhood on which the full nonconstant target star of the fixed finite
moment is finite.  We apply \eqref{eq:p21p3-gfc-balancing} to that entire
finite target star.  We never infer a balancing identity by deleting rays
and passing to an arbitrary finite substar.

\paragraph{Legitimate source-separated localization.}
We make explicit why no smooth plateau is being treated as an analytic
moment coordinate.  Let $\alpha$ be a compactly supported smooth
$(1,1)$-test form.  Start by localizing in the source.  The local moment
description \cite[Proposition~15.1.1(b)]{CLD25}, a smooth partition of
unity, compactness of $\operatorname{Supp}(\alpha)$, the finite-fibre
assertion \cite[Corollary~3.18]{GJR26}, and subdivision give finitely many
summands, each contained either in one relatively compact
virtual-annular collar or in one star neighborhood of a retained vertex.
The two collars at a type-$3$ point may be retained in its single star.
An edge neighborhood meets no other retained source edge, and a star
neighborhood contains no other retained vertex.  This source-first
operation is made before any Euclidean pushforward; distinct edge pieces
are never put into one compact source domain merely in order to separate
them by a cutoff.

Fix one summand in a source neighborhood $W$.  On an edge collar, apply
\cref{lem:p21p3-gf-legal-source-localization} directly.  At a type-$3$
center, first include in its prescribed block a $K$-analytic unit with
nonzero annular slope and shorten the two collars; the existence and
two-sided monomial norm formula are part of
\cite[\S\S4.2.11.2 and 4.3.9--4.3.10]{Ducros24CurveStructure}.  Apply the
same lemma after this shortening.

Suppose now that $W$ is a one-vertex star about a type-$2$ point $x$.
For the comparison at hand, let $A$ be the finite family of branches that
must be isolated; it contains the branches on which the form block is
nonconstant and, in the divisor calculation below, the finite support of
the principal divisor being tested.  We invoke
\cref{itm:p21p3-gf-bridge:localization} with $T$, this exact test summand
$\alpha$, this exact $W$, and the selected family $A$, and spell out its
two-stage output.  Its first application of
\cref{lem:p21p3-gf-legal-source-localization} produces a legal form and
retained-cutoff moment $F_0$.  After the branch block is appended, write
$G_0=(F_0,F_{\mathrm{br}})$.  The bridge then shrinks to its canonical-collar
one-vertex source neighborhood $W_1\subset W$.  Every virtual-disc
component is constant for $G_{0,\mathrm{trop}}$ on its entirety; only
finitely many retained virtual-annular components can be active, and every
selected active ray is separated from every other active branch.

Crucially, the bridge next reapplies
\cref{lem:p21p3-gf-legal-source-localization} \emph{inside $W_1$}, with
$G_0$ as its prescribed block.  We use its re-localized output
\begin{equation}
\label{eq:p21p3-gfc-relocalized-output}
 V\subset W_1,\qquad F,\qquad
 \Omega=\mathbf R^N\setminus
 F_{\mathrm{trop}}\bigl(\partial(V/W_1)\bigr),\qquad
 U=F_{\mathrm{trop}}^{-1}(\Omega).
\end{equation}
For every selected branch $b_i\in A$, it supplies the full canonical collar
$Z_{b_i}\subset U$ and the exact identity
\begin{equation}
\label{eq:p21p3-gfc-imported-exact-preimage}
 F_{\mathrm{trop}}^{-1}\!\left(
 F_{\mathrm{trop}}(x)+(0,\varepsilon_i)v_i\right)\cap U
 =Z_{b_i},
\end{equation}
which is precisely
\eqref{eq:p21p3-gf-source-separated-ray-preimage}.  Thus the isolated
target ray receives exactly its named source collar, not merely a branch
with the same ray germ.  This is the source-separation fact used below.

The full moment $F$ has finitely many analytic-unit coordinates.  By
\cite[Proposition~4.2.4(ii)]{Ducros24CurveStructure}, only finitely many
branches at $x$ are nonconstant for at least one of them.  Hence the
\emph{full} nonconstant target star in this localized moment is finite.
All of its rays, including rays activated by the second localization, are
retained when we apply Euclidean balancing.  We do not discard the
remaining geometric branches by a finite-fibre argument: they are constant
for $F$ on cofinal collars and map to the vertex.

In every case, the legal source-localization lemma gives a genuine analytic
block $H$ and a smooth $\Phi$ with
$\chi=\Phi\circ H_{\mathrm{trop}}$ whenever a source cutoff $\chi$ is
retained.  We append $H$, never $\chi$.  It also supplies the compact
domain $V$, the full moment $F$, and the proper map
$F_{\mathrm{trop}}|_U:U\to\Omega$; in the type-$2$ case these are exactly
the re-localized objects in
\eqref{eq:p21p3-gfc-relocalized-output}.  Thus no nonlinear cutoff is
declared to be a moment coordinate, and the current is never multiplied by
a cutoff.  The proper Euclidean pushforward
\begin{equation}
\label{eq:p21p3-gfc-euclidean-pushforward}
 \langle T_\Omega,\beta\rangle
       :=\langle T|_U,F^*\beta\rangle,
 \qquad \beta\in A_c^{1,1}(\Omega),
\end{equation}
defines a positive closed Euclidean current supported on a locally finite
one-dimensional polyhedral set.  Restriction to the open set $U$, not
multiplication by a cutoff, preserves closedness.  In a type-$2$ star,
\eqref{eq:p21p3-gfc-imported-exact-preimage} ensures that an isolated
selected ray receives the contribution of exactly its named source collar.

On such a collar, let $b_e\geq0$ be the flux of $T_\Omega$ from
\eqref{eq:p21p3-gfc-edge-density}, and let $m_e>0$ be the corresponding
flux for the canonical integration current.  Put
\begin{equation}
\label{eq:p21p3-gfc-branch-ratio}
                              a_e:=b_e/m_e.
\end{equation}
On a common source-separated concatenation of two moments, tropical
pushforward multiplies $b_e$ by the tangent-lattice expansion index:
the tangent tensor supplies its square and one-dimensional measure change
supplies its inverse.  The calibrated flux $m_e$ is multiplied by the
same index by \cite[\S14.3.2, Lemma~14.3.4, and
Proposition~14.4.2]{CLD25}.  Therefore $a_e$ is independent of the
moment, subdivision, and localized presentation, and is constant along
the virtual-annular collar.  This comparison is entirely downstairs over
$K$.

\paragraph{Comparison at type-$3$ and type-$2$ points.}
At an interior type-$3$ point the two branch rays in the refined moment
are opposite rays of one line.  Applying
\eqref{eq:p21p3-gfc-balancing} first to $T$ and then to $\delta_Y$
shows that the ratios \eqref{eq:p21p3-gfc-branch-ratio} on the two sides
are equal.  Possible virtual-annulus expansion factors occur equally in
the ordinary and calibrated fluxes and have already cancelled in the
ratio.

Let $x$ now be an interior type-$2$ point, and let $C_x$ be its
normal projective integral residual curve.  For a closed point
$p\in C_x$, use a rational function $s_p$ with
$\operatorname{ord}_p(s_p)=1$, a norm-one analytic lift, and
\cref{itm:p21p3-gf-bridge:ray}.  Apply
\cref{itm:p21p3-gf-bridge:localization} with the $p$-branch among the
selected branches.  Its re-localized exact-preimage identity
\eqref{eq:p21p3-gfc-imported-exact-preimage} makes the resulting ray the
image of precisely that source collar.  The outgoing downstairs vector
has a coordinate equal to one and is therefore primitive.  Denote the
calibrated flux on this isolated branch ray by $m_p>0$, and its ratio by
$a_p$.  These quantities are independent of the normalized moment: in
a common concatenation the source tangent and both projected tangents
retain a coordinate equal to one, so both projection indices are one.

Define additive maps on the divisor group of $C_x$ by
\begin{equation}
\label{eq:p21p3-gfc-divisor-functionals}
 L_T([p])=a_pm_p,
 \qquad
 L_Y([p])=m_p.
\end{equation}
They are nonnegative on effective divisors.  If
$g\in\operatorname{Rat}(C_x)^\times$, append a norm-one analytic lift
of $g$, and apply the source-separated localization in
\cref{itm:p21p3-gf-bridge:localization} simultaneously to the finite set
$\operatorname{Supp}(\operatorname{div}(g))$.  In the scalar coordinate
of this lift, the primitive slope on the $p$-branch is
$\operatorname{ord}_p(g)$; every branch outside that finite support has
slope zero.  The scalar component of
\eqref{eq:p21p3-gfc-balancing} gives
\begin{equation}
\label{eq:p21p3-gfc-T-principal}
       0=\sum_p a_pm_p\operatorname{ord}_p(g)
        =L_T(\operatorname{div}(g)).
\end{equation}
The integration current is closed by \cite[\S7.7.4]{CLD25}; applying the
same calculation to $\delta_Y$ gives
\begin{equation}
\label{eq:p21p3-gfc-Y-principal}
       0=\sum_p m_p\operatorname{ord}_p(g)
        =L_Y(\operatorname{div}(g)).
\end{equation}
No infinite sum and no cutoff on an infinite-valence star is used.

We use the following elementary consequence of Riemann--Roch over the
actual residue field.  Let $C$ be a normal projective integral curve over
a field $k$, and let
$L:\operatorname{Div}(C)\to\mathbf R$ be additive, nonnegative on
effective divisors, and zero on principal divisors.  If $D$ has degree
zero, choose an effective divisor $A$ such that
$\deg A+\chi(C,\mathcal O_C)>0$.
For every $n\in\mathbf Z$, Riemann--Roch gives
\[
 h^0(C,\mathcal O_C(A+nD))
 \geq \deg A+\chi(C,\mathcal O_C)>0.
\]
Thus $A+nD$ is linearly equivalent to an effective divisor and
\[
                         0\leq L(A)+nL(D)
                         \qquad(n\in\mathbf Z).
\]
Letting $n$ tend to both signs of infinity gives $L(D)=0$.  Choose
any closed point $p_0$ of $C$.  Since
$\deg(p_0)D-\deg(D)[p_0]$
has degree zero,
\begin{equation}
\label{eq:p21p3-gfc-degree-functional}
                 L(D)=c_L\deg(D),
 \qquad
                 c_L:=\frac{L([p_0])}{\deg(p_0)}\geq0.
\end{equation}
This argument neither assumes geometric integrality nor uses a degree-one
closed point.

Apply \eqref{eq:p21p3-gfc-degree-functional} to $L_Y$ and $L_T$.
There are $c_Y>0$ and $c_T\geq0$ such that
\[
                  m_p=c_Y\deg(p),
 \qquad
                  a_pm_p=c_T\deg(p)
\]
for every closed point $p$.  Hence
\begin{equation}
\label{eq:p21p3-gfc-equal-star-ratios}
                              a_p=c_T/c_Y
\end{equation}
on every branch at $x$.

\paragraph{Globalization on the smooth curve.}
Let $E\subset Y$ be the set of type-$1$ and type-$4$ points.  We justify
the endpoint assertion over the original field.  At a nonrigid point,
\cite[Theorem~4.2.5]{Ducros24CurveStructure} says that the existence of
two branches forces type $2$ or $3$; purity and the local graph structure
give at least one branch.  At a rigid point the virtual-disc neighborhood
of \cite[Theorem~4.5.4 and \S4.5.4.1]{Ducros24CurveStructure} gives the
same conclusion.  Thus every point of $E$ is unibranch, as is also
recorded for nonrigid type-$1$/type-$4$ points in
\cite[\S4.5.19.1]{Ducros24CurveStructure}.

Connected analytic spaces are arcwise connected by
\cite[\S12.1.12]{CLD25}.  If two points lie outside $E$, choose an arc
joining them.  By \cite[\S1.3.12]{Ducros24CurveStructure}, a point lying
in the open interval of a graph has at least the two branches supplied by
the two sides.  Hence no unibranch point of $E$ can occur in the interior
of the chosen arc, and $Y\setminus E$ is connected.  The ratio is
constant on each annular collar, propagates through type-$3$ points by
the two-ray comparison, and propagates through type-$2$ points by
\eqref{eq:p21p3-gfc-equal-star-ratios}.  It is therefore one constant
$c\geq0$ on $Y\setminus E$.

The source-separated local moments tropicalize every local test form.
The edge-density calculation and the absence of a vertex remainder then
show $T=c\delta_Y$ on a neighborhood of each point of $Y\setminus E$.
Every point of $E$ has tropical dimension zero, so
\cite[Proposition~3.3.6]{CLD25} says that the germ of every smooth
$(1,1)$-test form vanishes there.  Thus the support of each such test
form is disjoint from $E$, and the current $T-c\delta_Y$, already zero
off $E$, pairs to zero with every test form.  This proves
\eqref{eq:p21p3-gfc-smooth-rigidity}.  Uniqueness follows by testing on a
nonempty virtual-annular band.  All arguments were made on compact test
supports, so $Y$ need not be compact.
\end{proof}

\begin{proof}[Proof of \cref{prop:p21p3-gfc-curve-rigidity}]
We first pass from the smooth normalizations to the reduced curve.  On a
good affinoid domain write
\[
 \mathcal Q:=\nu_*\mathcal O_{\widetilde X}/\mathcal O_X,
 \qquad
 \mathfrak c:=\operatorname{Ann}_{\mathcal O_X}(\mathcal Q),
\]
and let $Z$ be the union of
\[
 \operatorname{Supp}(\mathcal Q)
   =V(\mathfrak c)
   =\operatorname{Supp}(\mathcal O_X/\mathfrak c)
\]
with the pairwise intersections of distinct irreducible components.
Analytic local rings on a good space are excellent by
\cite[\S11.1.10]{CLD25}; normalization is finite and its components
correspond to the irreducible components by
\cite[\S11.3.1]{CLD25}.  Hence $\mathcal Q$ and the conductor are
coherent.  Their support misses every generic point of the reduced pure
curve and is therefore zero-dimensional.  The same dimension argument
applies to proper intersections of distinct components.  Irreducible
components are locally finite by \cite[\S11.2.5]{CLD25}; consequently
$Z$ is closed and locally finite.

Strictness is used here.  A zero-dimensional strict affinoid algebra is
Artinian, and the strict affinoid Nullstellensatz makes the completed
residue field of each of its points finite over $K$.  Thus every point
of $Z$ is rigid.  A finite morphism induces a finite extension of
completed residue fields, so every point of
\begin{equation}
\label{eq:p21p3-gfc-conductor-preimage}
                         Z_i:=\nu_i^{-1}(Z),
\end{equation}
is rigid as well, though it need not be $K$-rational.  Moreover, a finite
algebra over a strict affinoid algebra is strict.  Since normalization is
finite, every $\widetilde X_i$ is therefore strict.

The finite normalization map is topologically proper, has empty relative
boundary, and is an isomorphism away from $Z_i$.  Applying the
boundary-composition identity \cite[(11.4.1.3)]{CLD25} to
$X_i\to X\to\mathcal M(K)$ and then to
$\widetilde X_i\to X_i\to\mathcal M(K)$ shows that
$\widetilde X_i$ is boundaryless.  It is connected because it is the
normalization of one irreducible component.

The fundamental-cycle identity is
\begin{equation}
\label{eq:p21p3-gfc-fundamental-cycle}
             \delta_X=\sum_i(\nu_i)_*\delta_{\widetilde X_i}.
\end{equation}
Indeed, a reduced pure curve has generic multiplicity one on every
irreducible component.  Integration agrees with the associated cycle by
\cite[Proposition~16.6.3]{CLD25}, and normalization has calibrated degree
one over the normal locus.  The complement is the locally finite rigid
set $Z$, where every $(1,1)$-test-form germ vanishes by
\cite[Proposition~3.3.6]{CLD25}; this proves
\eqref{eq:p21p3-gfc-fundamental-cycle} everywhere.

The complement $\widetilde X_i\setminus Z_i$ is open, good, separated,
strict, rig-smooth, and boundaryless.  Deleting $Z_i$ does not disconnect
$\widetilde X_i$.  Indeed, two
remaining points can be joined by an arc, and
\cite[\S1.3.12]{Ducros24CurveStructure} shows that no rigid unibranch
point can lie in the interior of that arc.  On $X\setminus Z$, the component
images are disjoint and normalization is an analytic isomorphism.
Applying \cref{prop:p21p3-gfc-smooth-rigidity} to each connected
$\widetilde X_i\setminus Z_i$ gives
\[
                         T=\sum_i c_i\delta_{X_i},
                         \qquad c_i\geq0,
\]
there.  Set $S=T-\sum_i c_i\delta_{X_i}$.  It has zero germ off $Z$,
while every $(1,1)$-test form has support disjoint from $Z$ by
\cite[Proposition~3.3.6]{CLD25}.  Hence $S=0$, proving
\eqref{eq:p21p3-gfc-curve-classification}.  The sum is locally finite,
and uniqueness follows by restricting to an open normal locus of $X_i$
meeting no other component.

It remains to prove the Bedford--Taylor assertion.  Put
$\widetilde W_i=\nu_i^{-1}(W)$.  Multiplying the order-zero identity
\eqref{eq:p21p3-gfc-fundamental-cycle} by a continuous function $F$ on
$W$, using \cite[Example~7.2.12]{CLD25}, gives
\begin{equation}
\label{eq:p21p3-gfc-function-cycle-decomposition}
 [F]=F\delta_W
     =\sum_i(\nu_i)_*
        \bigl((\nu_i^*F)\delta_{\widetilde W_i}\bigr).
\end{equation}
The projection identity used here follows directly by testing against a
smooth form, and the sum is finite on each test support.

Pullback preserves continuous plurisubharmonic functions by
\cite[Proposition~8.1.5]{CLD25}.  On $\widetilde W_i$, set
\[
 \mu_i^U=\mathrm d'\mathrm d''[\nu_i^*U],
 \qquad
 \mu_i^V=\mathrm d'\mathrm d''[\nu_i^*V].
\]
These are positive Radon measures by \cite[Lemma~8.5.5]{CLD25}.
Proper pushforward commutes with the differential on currents by
transposition; equivalently, use \cite[Proposition~8.5.6]{CLD25}.
Applying \eqref{eq:p21p3-gfc-function-cycle-decomposition} to $U$ and
$V$, the curvature equality
\eqref{eq:p21p3-gfc-curve-equal-curvature} becomes
\begin{equation}
\label{eq:p21p3-gfc-component-measure-sum}
                    \sum_i(\nu_i)_*(\mu_i^U-\mu_i^V)=0.
\end{equation}

Away from $Z$, the component images are disjoint and $\nu_i$ is an
isomorphism.  Hence each signed measure $\mu_i^U-\mu_i^V$ is supported
on the locally finite rigid set $Z_i\cap\widetilde W_i$.  At any point of
this set, apply \cref{lem:p21p3-gfc-rigid-zero-atom} on a connected open
neighborhood contained in $\widetilde W_i$.  Neither positive measure has
an atom there.  A signed Radon measure supported on a locally finite set
is the locally finite sum of its point masses.  Therefore
\begin{equation}
\label{eq:p21p3-gfc-component-measure-equality}
                              \mu_i^U=\mu_i^V
                              \quad\text{for every }i.
\end{equation}

Finally put $h=U-V$.  By
\eqref{eq:p21p3-gfc-curve-classification},
\[
 hT|_W=\sum_i c_i(\nu_i)_*
       \bigl((\nu_i^*h)\delta_{\widetilde W_i}\bigr).
\]
Proper pushforward commutes with $\mathrm d'\mathrm d''$, and
\eqref{eq:p21p3-gfc-component-measure-equality} gives
\begin{align*}
 \mathrm d'\mathrm d''(hT|_W)
 &=\sum_i c_i(\nu_i)_*
       \mathrm d'\mathrm d''[\nu_i^*U-\nu_i^*V] \\
 &=\sum_i c_i(\nu_i)_*(\mu_i^U-\mu_i^V)=0.
\end{align*}
This is \eqref{eq:p21p3-gfc-curve-bt-conclusion}.  It proves the
independence condition of \cite[Definition~8.6.2]{CLD25} for the
continuous psh pairs stated here; no claim about discontinuous psh
functions is made.
\end{proof}


% --- Fresh conditional nonreduced addendum proof begins here. ---

% Fresh source-aware proof of the canonical-collar nonreduced addendum.
% Date: 28 August 2026.
% This proof and its exact reduced prefix are not integrated into the book.
% They remain conditional under the two-stage production firewall stated
% in \cref{rem:p21p3-gfnr-scope}.

\begin{proof}
The reduction has the same underlying topological space and the same
completed residue fields as \(X\).  Affinoid charts pass to quotients by
nilpotent ideals, so goodness, separatedness, strictness, and pure
dimension one pass to \(X_{\mathrm{red}}\).  The immersion \(\iota\) is
finite and hence has empty relative boundary by
\cite[\S11.4.4]{CLD25}.  The boundary-composition identity
\cite[(11.4.1.3)]{CLD25} therefore shows that
\(X_{\mathrm{red}}\) is boundaryless.  Thus \(X_{\mathrm{red}}\)
satisfies the ambient hypotheses of the conditional reduced theorem
\cref{prop:p21p3-gfc-curve-rigidity}.

\smallskip
\emph{Forms, currents, and psh functions.}
Because \(X\) is good, it is almost good.  The reduction-invariance result
\cite[Proposition~15.1.3(b) and the following form variant]{CLD25}
therefore gives an isomorphism

\begin{equation}
\label{eq:p21p3-gfnr-form-isomorphism}
       \iota^*:A_X^{p,q}\xrightarrow{\ \sim\ }
                       A_{X_{\mathrm{red}}}^{p,q}.
\end{equation}

The point needed for currents is topological, not merely algebraic.  By
\cite[\S\S~6.4.1 and 7.1.1--7.1.2]{CLD25}, a datum for the compact-test
topology consists of a compact support and a finite family of moment charts.
Restriction sends such a datum on \(X\)
to the same datum on \(X_{\mathrm{red}}\).  Conversely, the cited
almost-good result applies because \(X\) is good.  At every point of the
compact support, a finite tuple of invertible functions defining a reduced
moment chart lifts through the nilpotent quotient to a tuple of unit germs
on \(X\).  Each lift has the same absolute value as its reduction at every
analytic point.  After shrinking the ordinary open neighbourhoods and
using compactness, one obtains a finite lifted refinement; product moment
maps give a common refinement on overlaps.

Restriction and these finite lifted refinements are cofinal systems of
compact chart data for the final locally convex topologies.  On a lifted
chart, the tropical images, Lagerberg representatives, polyhedral
refinements, and every \(C^r\)-seminorm are identical.  Consequently

\begin{equation}
\label{eq:p21p3-gfnr-test-homeomorphism}
 \iota^*:A^{p,q}_{X,c}\xrightarrow{\ \sim\ }
                         A^{p,q}_{X_{\mathrm{red}},c}
\end{equation}

is a homeomorphism.  The same chart description proves compatibility with
wedge product, \(J\), the differentials, supports, and all positivity cones
of forms.  Taking continuous duals gives

\begin{equation}
\label{eq:p21p3-gfnr-current-isomorphism}
 \iota_*:D^{a,b}(X_{\mathrm{red}})
              \xrightarrow{\ \sim\ }D^{a,b}(X),
 \qquad
 \langle\iota_*S,\alpha\rangle
       =\langle S,\iota^*\alpha\rangle .
\end{equation}

This isomorphism identifies the weak topologies.  It preserves positivity
because the test-form cones agree, and it preserves closedness because
\(\iota^*\) commutes with \(\mathrm d'\) and \(\mathrm d''\).
Multiplication by a smooth function commutes with \(\iota_*\) directly by
transposition.  If \(S\) is positive, the Borel product of
\cite[Example~7.2.12]{CLD25} is defined through the measures
\(\mu_{S,\omega}\) attached to positive test forms \(\omega\).  Under
\(\iota^*\), the positive test forms and these attached measures agree.
Thus \(\iota_*(fS)=f\iota_*S\) for every locally bounded Borel function
\(f\).

If \(S=S_1-S_2\) locally with \(S_1,S_2\) positive, define
\(fS=fS_1-fS_2\).  This is independent of the decomposition.  Indeed, from
\(S_1-S_2=S'_1-S'_2\) one gets
\(S_1+S'_2=S'_1+S_2\); equality of these positive currents gives equality
of all their attached measures, and hence
\(fS_1+fS'_2=fS'_1+fS_2\).  The local definitions glue, and the same
calculation proves compatibility with \(\iota_*\).  Nothing in this
argument defines \(fS\) for an arbitrary order-zero current.

For the psh assertion, stage zero of the CLD filtration is the same on the
two spaces: smoothness and positivity of
\(\mathrm d'\mathrm d''u\) agree by
\eqref{eq:p21p3-gfnr-form-isomorphism}.  Successor stages take decreasing
limits, the continuous filtration takes local uniform limits, and limit
ordinal stages take local unions.  The underlying open sets and functions
are identical.  Transfinite induction in the filtration of
\cite[Lemma~8.1.4 and Proposition~8.1.11]{CLD25} therefore identifies the
psh and continuous psh sheaves at every stage and after stabilization.  The
\(G\)-form argument is the same, using the \(G\)-form variant immediately after
\cite[Proposition~15.1.3]{CLD25}.  This proves
\cref{itm:p21p3-gfnr:calculus}.

\smallskip
\emph{The integration current and generic multiplicity.}
By \cite[Definition~13.2.5 and \S13.5.2]{CLD25}, the fundamental cycles
are

\[
                 [X]=\sum_i m_i[X_i],
                 \qquad
                 [X_{\mathrm{red}}]=\sum_i[X_i].
\]

Let
\(\bar\nu_i:\widetilde X_i\to X_i\hookrightarrow X_{\mathrm{red}}\),
so that \(\nu_i=\iota\circ\bar\nu_i\).  The normalization is finite and
generically an isomorphism.  The compatibility of integration with analytic
cycles in \cite[Proposition~16.6.3]{CLD25} identifies the current of
\([X_i]\) with
\((\bar\nu_i)_*\delta_{\widetilde X_i}\) on \(X_{\mathrm{red}}\), and
with \((\nu_i)_*\delta_{\widetilde X_i}\) after pushforward to \(X\).
It follows that

\[
 \delta_X=\sum_i m_i\delta_i,
 \qquad
 \iota_*\delta_{X_{\mathrm{red}}}=\sum_i\delta_i.
\]

The irreducible components are locally finite by
\cite[\S11.2.5]{CLD25}, so both sums are locally finite.  Their equality
forces \(m_i=1\) for every \(i\): restrict to a nonempty normal open subset
of \(X_i\) meeting no other component and use the nonzero integration
current there.  Conversely, \(m_i=1\) for every \(i\) visibly gives
equality.  Since the length of the generic local ring along \(X_i\) is
\(m_i\), this is equivalent to generic reducedness along every component.
This proves \cref{itm:p21p3-gfnr:integration}.

\smallskip
\emph{Classification.}
Let \(T\) be positive and closed on \(X\).  By
\eqref{eq:p21p3-gfnr-current-isomorphism}, there is a unique positive
closed current \(T_{\mathrm{red}}\) on \(X_{\mathrm{red}}\) with
\(T=\iota_*T_{\mathrm{red}}\).  The conditional reduced classification
\cref{itm:p21p3-gfc-curve:classification} of
\cref{prop:p21p3-gfc-curve-rigidity} gives

\[
                  T_{\mathrm{red}}=\sum_i c_i
                    (\bar\nu_i)_*\delta_{\widetilde X_i},
                  \qquad c_i\geq0.
\]

Pushing forward by \(\iota\) proves
\eqref{eq:p21p3-gfnr-classification}, including local finiteness and
uniqueness.  If \(X_{\mathrm{red}}\) is connected and rig-smooth, its
regular local rings are domains, so its irreducible components are
disjoint open-and-closed subsets.  It therefore has one component, with
one generic multiplicity \(m\).  Then
\(\delta_X=m\delta_1\), and
\(c_1\delta_1=(c_1/m)\delta_X\).  This proves
\cref{itm:p21p3-gfnr:classification}.

\smallskip
\emph{The strict-rigid exceptional locus.}
Let
\(\bar\nu:\coprod_i\widetilde X_i\to X_{\mathrm{red}}\) be the
normalization map.  On a good affinoid domain put

\[
 \mathcal Q:=\bar\nu_*\mathcal O_{\coprod_i\widetilde X_i}
                 /\mathcal O_{X_{\mathrm{red}}},
 \qquad
 \mathfrak c:=\operatorname{Ann}_{\mathcal O_{X_{\mathrm{red}}}}
                         (\mathcal Q).
\]

Let \(Z\) be the union of the conductor locus

\[
 \operatorname{Supp}(\mathcal Q)
   =V(\mathfrak c)
   =\operatorname{Supp}(\mathcal O_{X_{\mathrm{red}}}/\mathfrak c)
\]

and the pairwise intersections of distinct irreducible components.
Analytic local rings on a good space are excellent by
\cite[\S11.1.10]{CLD25}; normalization is finite and its components
correspond to the irreducible components by
\cite[\S11.3.1]{CLD25}.  Hence \(\mathcal Q\) and the conductor are
coherent.  Their support misses every generic point of the reduced pure
curve and is zero-dimensional.  The same dimension argument applies to
the proper intersections of distinct components.  Local finiteness of the
components, again \cite[\S11.2.5]{CLD25}, shows that \(Z\) is closed and
locally finite.

Here strictness replaces algebraic closedness.  Locally, a
zero-dimensional closed analytic subspace of a strict affinoid domain is
the spectrum of a zero-dimensional strict affinoid algebra \(B\).  By
\cite[Theorem~4.5]{Tate71Rigid}, the algebra \(B\) is noetherian; having
Krull dimension zero, it is therefore Artinian.
For a point \(z\) of its Berkovich spectrum, the kernel of the associated
seminorm is maximal; the strict affinoid Nullstellensatz
\cite[Theorem~4.5]{Tate71Rigid} makes
\(B/\ker(z)\) finite over \(K\).  This finite-dimensional \(K\)-algebra is
already complete, so
\(\mathcal H(z)=B/\ker(z)\) is finite over \(K\).  Thus every point of
\(Z\) is rigid, although it need not be \(K\)-rational.

Put

\begin{equation}
\label{eq:p21p3-gfnr-conductor-preimage}
                         Z_i:=\bar\nu_i^{-1}(Z).
\end{equation}

A finite morphism induces finite extensions of completed residue fields,
so every point of \(Z_i\) is rigid.  A finite algebra over a strict
affinoid algebra is strict by \cite[Theorem~4.4]{Tate71Rigid};
consequently each \(\widetilde X_i\) is strict.  The finite maps
\(X_i\to X_{\mathrm{red}}\) and
\(\widetilde X_i\to X_i\) have empty relative boundary by
\cite[\S11.4.4]{CLD25}.  Applying
\cite[(11.4.1.3)]{CLD25} twice shows that every
\(\widetilde X_i\) is boundaryless.  It is good and separated, and it is
rig-smooth by hypothesis.  Away from \(Z\), the component images are
disjoint and each \(\bar\nu_i\) is an analytic isomorphism.  These are
exactly the hypotheses needed below for the rigid zero-atom input; no
assertion about arbitrary type-\(1\) points is used.

\smallskip
\emph{The multiplicity-weighted curvature identity.}
Fix \(W,U,V\) as in \cref{itm:p21p3-gfnr:bedford-taylor}, identify \(W\)
with the corresponding open subset of \(X_{\mathrm{red}}\), and put
\(\widetilde W_i=\bar\nu_i^{-1}(W)\).  Denote the base-changed maps by

\[
 \nu_{i,W}^{X}:\widetilde W_i\longrightarrow W\subset X,
 \qquad
 \nu_{i,W}^{\mathrm{red}}:\widetilde W_i\longrightarrow
                         W\subset X_{\mathrm{red}}.
\]

Pullback preserves continuous psh functions by
\cite[Proposition~8.1.5]{CLD25}.  On \(\widetilde W_i\), define positive
Radon measures

\[
 \mu_i^U:=\mathrm d'\mathrm d''[\bar\nu_i^*U],
 \qquad
 \mu_i^V:=\mathrm d'\mathrm d''[\bar\nu_i^*V],
 \qquad
 \sigma_i:=\mu_i^U-\mu_i^V.
\]

Their positivity is \cite[Lemma~8.5.5]{CLD25}.  For a continuous function
\(F\) on \(W\), multiply the two locally finite identities in
\eqref{eq:p21p3-gfnr-fundamental-cycle} by \(F\), using
\cite[Example~7.2.12]{CLD25}.  Evaluation on a compactly supported test
form and the defining transpose pairing for finite pushforward give

\begin{align}
 [F]_X
   &=\sum_i m_i(\nu_{i,W}^{X})_*
       \bigl(((\nu_{i,W}^{X})^*F)\delta_{\widetilde W_i}\bigr),
 \label{eq:p21p3-gfnr-function-current-X}\\
 [F]_{X_{\mathrm{red}}}
   &=\sum_i(\nu_{i,W}^{\mathrm{red}})_*
       \bigl(((\nu_{i,W}^{\mathrm{red}})^*F)
                    \delta_{\widetilde W_i}\bigr).
 \label{eq:p21p3-gfnr-function-current-red}
\end{align}

For \(F=U,V\), differentiate these identities and use the continuous-psh
projection formula \cite[Proposition~8.5.6]{CLD25}; equivalently, proper
pushforward commutes with the current differentials by transposition.
Thus equality of the curvature currents computed with \(\delta_X\) is
equivalent to

\begin{equation}
\label{eq:p21p3-gfnr-weighted-measure-sum}
                       \sum_i m_i(\nu_{i,W}^{X})_*\sigma_i=0,
\end{equation}

whereas equality computed with \(\delta_{X_{\mathrm{red}}}\) is
equivalent to

\begin{equation}
\label{eq:p21p3-gfnr-unweighted-measure-sum}
                       \sum_i
                         (\nu_{i,W}^{\mathrm{red}})_*\sigma_i=0.
\end{equation}

Write \(Z_W=Z\cap W\) and
\(Z_{i,W}=Z_i\cap\widetilde W_i=\bar\nu_i^{-1}(Z_W)\).  Away from \(Z_W\),
the component images are disjoint and the relevant normalization map is
an isomorphism.  Restricting either
\eqref{eq:p21p3-gfnr-weighted-measure-sum} or
\eqref{eq:p21p3-gfnr-unweighted-measure-sum} to such an open set, and
using \(m_i>0\) in the weighted case, yields

\begin{equation}
\label{eq:p21p3-gfnr-component-support}
                       \operatorname{Supp}(\sigma_i)
                       \subseteq Z_{i,W}
                       \quad\text{for every }i.
\end{equation}

Every point of \(Z_{i,W}\) is rigid.  Around such a point, choose an open
neighbourhood contained in \(\widetilde W_i\).  It is a good separated
strict rig-smooth boundaryless curve, so
\cref{lem:p21p3-gfc-rigid-zero-atom} applies to the restrictions of
\(\bar\nu_i^*U\) and \(\bar\nu_i^*V\).  Neither positive measure
\(\mu_i^U\) nor \(\mu_i^V\) has an atom there.  A signed Radon measure
supported on a locally finite set is the locally finite sum of its point
masses.  Hence \eqref{eq:p21p3-gfnr-component-support} implies

\begin{equation}
\label{eq:p21p3-gfnr-component-measure-equality}
                         \sigma_i=0\qquad\text{for every }i.
\end{equation}

This argument works starting from either the weighted or the unweighted
identity.  Conversely,
\eqref{eq:p21p3-gfnr-component-measure-equality} implies both identities.
This proves the exact curvature equivalence
\eqref{eq:p21p3-gfnr-curvature-equivalence}; in particular, generic
multiplicities cannot conceal cancellation at a component intersection.

Finally put \(h=U-V\) and write \(T=\sum_i c_i\delta_i\).  Multiplication
by \(h\) is defined for the positive current \(T\) by
\cite[Example~7.2.12]{CLD25}; continuity makes \(h\) locally bounded.
Proper pushforward, linearity of the current differentials, and
\eqref{eq:p21p3-gfnr-component-measure-equality} give

\begin{align*}
 \mathrm d'\mathrm d''(hT|_W)
 &=\sum_i c_i(\nu_{i,W}^{X})_*
     \mathrm d'\mathrm d''
       [\bar\nu_i^*U-\bar\nu_i^*V]\\
 &=\sum_i c_i(\nu_{i,W}^{X})_*(\mu_i^U-\mu_i^V)=0.
\end{align*}

This is \eqref{eq:p21p3-gfnr-bt-conclusion}.  It proves precisely the
Bedford--Taylor independence identity for the continuous psh pairs stated
in the proposition.  The proof does not invoke, and the proposition does
not assert, the additional discontinuous-psh quantifier in
\cite[Definition~8.6.2]{CLD25}.  This completes the proof.
\end{proof}
